\documentclass[10pt,a4paper]{amsart}
\usepackage[margin=19mm]{geometry}
\usepackage{amsmath,amssymb,mathtools}
\usepackage[scr=rsfs,cal=euler]{mathalfa}
\usepackage{xfrac}
\usepackage[unicode,hidelinks,bookmarks=false]{hyperref}
\newcommand{\ie}{i.e.\ }
\newcommand{\cf}{cf.\ }
\newcommand{\resp}{respectively\ }
\newcommand{\Subsection}{\subsection}
\providecommand{\nobreakdash}{\nobreak\hbox{-}}
\setlength{\emergencystretch}{2em}
\multlinegap0pt
\usepackage{amssymb}
\usepackage{dsfont}
\newtheorem{corollary}{Corollary}
\newtheorem{proposition}{Proposition}
\newcommand{\I}{\mathds{1}}
\newcommand{\NN}{\mathbb{N}}
\newcommand{\TT}{\mathbb{T}}
\newcommand{\ZZ}{\mathbb{Z}}
\newcommand{\cA}{\mathscr{A}}
\newcommand{\comp}{\circ}
\title{Braided quantum $\mathrm{SU}(2)$ group -- a case study}
\author{Jacek Krajczok, Piotr. M. So{\l}tan}
\date{Source: arXiv:2604.14937v1 (CC BY 4.0)}
\begin{document}
\maketitle

\noindent\emph{Source and licence.} Braided quantum $\mathrm{SU}(2)$ group -- a case study. Version arXiv:2604.14937v1 (CC BY 4.0), distributed by arXiv under the Creative Commons Attribution 4.0 International licence, http://creativecommons.org/licenses/by/4.0/, as recorded in the arXiv licence metadata for this exact version and shown on the abstract page https://arxiv.org/abs/2604.14937v1. Full licence notice: \texttt{licenses/arxiv-2604.14937-CC-BY-4.0.txt}.\par\smallskip

\noindent\textit{Editorial scope.} Counted unit: the article's proposition prop-lemma12 (source0.tex lines 724--729). The article's complete original proof of it is delivered with it (source0.tex lines 731--845, one \texttt{proof} environment of 115 lines). No mathematics is written, repaired or re-derived: the statement and the proof are the article's own lines, in the article's own notation, and no retained line is substituted or re-flowed.  Retained uncounted as the source-local context the proof needs: lines 267--273.  Nothing was omitted from any retained span, so the package carries no omission record.  No retained line contains a citation, so the wrapper carries no bibliography and no bibliography entry.  No retained line contains a figure environment, an image-inclusion command or drawing code, so no image is delivered and the package declares no asset dependency.  Editorial formatting only, in addition: an AMS \texttt{amsart} preamble that restates the article's own declarations and macros used by the retained lines, namely the environments \texttt{corollary}, \texttt{proposition} on separate counters, together with the standard packages those lines need.  The retained environments print in this document as Corollary 1, Proposition 1 in order of appearance, whereas the article prints the same items with the numbering of its own full document.  The complete native source of the article as distributed in the arXiv e-print for this exact version is bundled unchanged as \texttt{sources/198567/source0.tex}.
\begin{corollary}\label{cor-bazacA}
The system
\begin{equation}\label{eq-bazacA}
\bigl\{\alpha^n\gamma^m{\gamma^*}^k\,\bigr|\bigl.\,n,m,k\in\ZZ_+\bigr\}\cup\bigl\{{\alpha^*}^n\gamma^m{\gamma^*}^k\,\bigr|\bigl.\,n\in\NN,\:m,k\in\ZZ_+\bigr\}
\end{equation}
is a basis of $\cA$.
\end{corollary}

\begin{proposition}\label{prop-lemma12}
The antipode satisfies
\begin{equation}\label{eq-eq45}
\Delta_A\comp S=(S\boxtimes S)\comp\chi^{\boxtimes}\comp\Delta_A.
\end{equation}
\end{proposition}

\begin{proof}
First we check that both sides of \eqref{eq-eq45} coincide on the elements $\I$ (trivial), $\alpha$, $\alpha^*$, $\gamma$ and $\gamma^*$. We have
\begin{align*}
\Delta_A\bigl(S(\alpha)\bigr)
&=\Delta_A(\alpha^*)
=\bigl(\iota_1(\alpha)\iota_2(\alpha)-q\iota_1(\gamma)^* \iota_2(\gamma)\bigr)^*
=\iota_2(\alpha^*)\iota_1(\alpha^*)-\overline{q}\iota_2(\gamma^*)\iota_1(\gamma)\\
&=\iota_1(\alpha^*)\iota_2(\alpha^*)-\overline{q}\zeta\iota_1(\gamma)\iota_2(\gamma^*)
=\iota_1(\alpha^*)\iota_2(\alpha^*)-q\iota_1(\gamma)\iota_2(\gamma^*)
\end{align*}
and
\begin{align*}
(S\boxtimes S)\chi^{\boxtimes}\Delta_A(\alpha)
&=(S\boxtimes S)\chi^{\boxtimes}(\iota_1(\alpha)\iota_2(\alpha)-q\iota_1(\gamma^*)\iota_2(\gamma))\\
&=(S\boxtimes S)(\iota_1(\alpha)\iota_2(\alpha)-q\zeta\iota_1(\gamma)\iota_2(\gamma^*))\\
&=\iota_1(\alpha^*)\iota_2(\alpha^*)-q\zeta\iota_1(-\overline{q}\gamma)\iota_2(-q^{-1}\gamma^*)
=\iota_1(\alpha^*)\iota_2(\alpha^*)-q\iota_1(\gamma)\iota_2(\gamma^*).
\end{align*}
Similarly $\Delta_A\bigl(S(\alpha^*)\bigr)
=\Delta_A(\alpha)
=\iota_1(\alpha)\iota_2(\alpha)-q\iota_1(\gamma)^*\iota_2(\gamma)$ and
\begin{align*}
(S\boxtimes S)\chi^{\boxtimes}\Delta_A(\alpha^*)
&=(S\boxtimes S)\chi^{\boxtimes}\bigl((\iota_1(\alpha)\iota_2(\alpha)-q\iota_1(\gamma^*)\iota_2(\gamma))^*\bigr)\\
&=(S\boxtimes S)\chi^{\boxtimes}\bigl(\iota_2(\alpha^*)\iota_1(\alpha^*)-\overline{q}\iota_2(\gamma^*)\iota_1(\gamma)\bigr)\\
&=(S\boxtimes S)\chi^{\boxtimes}\bigl(\iota_1(\alpha^*)\iota_2(\alpha^*)-\overline{q}\zeta\iota_1(\gamma)\iota_2(\gamma^*)\bigr)\\
&=(S\boxtimes S)\bigl(\iota_1(\alpha^*)\iota_2(\alpha^*)-q\zeta\iota_1(\gamma^*)\iota_2(\gamma)\bigr)\\
&=\iota_1(\alpha)\iota_2(\alpha)-q\zeta\iota_1(-q^{-1}\gamma^*)\iota_2(-\overline{q}\gamma)\\
&=\iota_1(\alpha)\iota_2(\alpha)-q\iota_1(\gamma^*)\iota_2(\gamma),
\end{align*}
as well as $\Delta_A\bigl(S(\gamma)\bigr)
=-\overline{q}\Delta_A(\gamma)=-\overline{q}\iota_1(\gamma)\iota_2(\alpha)
-\overline{q}\iota_1(\alpha^*)\iota_2(\gamma)$ and
\begin{align*}
(S\boxtimes S)\chi^{\boxtimes}\Delta_A(\gamma)
&=(S\boxtimes S)\chi^{\boxtimes}\bigl(\iota_1(\gamma)\iota_2(\alpha)+\iota_1(\alpha^*)\iota_2(\gamma)\bigr)\\
&=(S\boxtimes S)\bigl(\iota_1(\alpha)\iota_2(\gamma)+\iota_1(\gamma)\iota_2(\alpha^*)\bigr)\\
&=\iota_1(\alpha^*)\iota_2(-\overline{q}\gamma)+\iota_1(-\overline{q}\gamma)\iota_2(\alpha)
\end{align*}
and finally
\begin{align*}
\Delta_A\bigl(S(\gamma^*)\bigr)
&=-q^{-1}\Delta_A(\gamma^*)
=-q^{-1}\bigl(\iota_1(\gamma)\iota_2(\alpha)+\iota_1(\alpha^*)\iota_2(\gamma)\bigr)^*\\
&=-q^{-1}\iota_2(\alpha^*)\iota_1(\gamma^*)-q^{-1}\iota_2(\gamma^*)\iota_1(\alpha)\\
&=-q^{-1}\iota_1(\gamma^*)\iota_2(\alpha^*)-q^{-1} \iota_1(\alpha)\iota_2(\gamma^*)
\end{align*}
and
\begin{align*}
(S\boxtimes S)\chi^{\boxtimes}\Delta_A(\gamma^*)
&=(S\boxtimes S)\chi^{\boxtimes}\bigl(\iota_1(\gamma^*)\iota_2(\alpha^*)+\iota_1(\alpha)\iota_2(\gamma^*)\bigr)\\
&=(S\boxtimes S)\bigl(\iota_1(\alpha^*)\iota_2(\gamma^*)+\iota_1(\gamma^*)\iota_2(\alpha)\bigr)\\
&=\iota_1(\alpha)\iota_2(-q^{-1}\gamma^*)+\iota_1(-q^{-1}\gamma^*)\iota_2(\alpha^*).
\end{align*}

Now let us take $x,y\in\cA$ and assume that both sides of \eqref{eq-eq45} have equal value for $x$ and $y$. Write (in Sweedler notation) $\Delta_A(x)=\sum \iota_1(x_{(1)})\iota_2(x_{(2)})$ assuming additionally that each $x_{(1)}$, $x_{(2)}$ is homogeneous, and similarly for $y$. Note that since $\Delta_A$ is $\TT$-equivariant, we can additionally assume $\deg(x)=\deg(x_{(1)})+\deg(x_{(2)})$ for each summand. Then writing $\lambda_{p,q}$ for the numerical factor $\zeta^{-\deg(p)\deg(q)}$ we compute
\begin{align*}
\Delta_A\bigl(S(xy)\bigr)
&=\lambda_{x,y}\Delta_A\bigl(S(y)S(x)\bigr)
=\lambda_{x,y}\Delta_A\bigl(S(y)\bigr)\Delta_A\bigl(S(x)\bigr)\\
&=\lambda_{x,y}\bigl((S\boxtimes S)\chi^\boxtimes\Delta_A(y)\bigr)
\bigl((S\boxtimes S)\chi^\boxtimes\Delta_A(x)\bigr)\\
&=\lambda_{x,y}\sum
\lambda_{y_{(1)},y_{(2)}}
(S\boxtimes S)\bigl(\iota_1(y_{(2)})\iota_2(y_{(1)})\bigr)
\lambda_{x_{(1)},x_{(2)}}(S\boxtimes S)\bigl(\iota_1(x_{(2)})\iota_2(x_{(1)})\bigr)\\
&=\lambda_{x,y}
\sum\lambda_{y_{(1)},y_{(2)}}
\lambda_{x_{(1)},x_{(2)}}
\iota_1\bigl(S(y_{(2)})\bigr)\iota_2\bigl(S(y_{(1)})\bigr)
\iota_1\bigl(S(x_{(2)})\bigr)\iota_2\bigl(S(x_{(1)})\bigr)\\
&=\lambda_{x,y}
\sum\lambda_{y_{(1)},y_{(2)}}
\lambda_{x_{(1)},x_{(2)}}\lambda_{x_{(2)},y_{(1)}}\iota_1\bigl(S(y_{(2)})\bigr)
\iota_1\bigl(S(x_{(2)})\bigr)\iota_2\bigl(S(y_{(1)})\bigr)\iota_2\bigl(S(x_{(1)})\bigr)
\end{align*}
and
\begin{align*}
(S\boxtimes S)&\chi^\boxtimes\Delta_A(xy)
=(S\boxtimes S)\chi^\boxtimes\bigl(\Delta_A(x)\Delta_A(y)\bigr)\\
&=\sum(S\boxtimes S)\chi^\boxtimes
\bigl(\iota_1(x_{(1)})\iota_2(x_{(2)})\iota_1(y_{(1)})\iota_2(y_{(2)})\bigr)\\
&=\sum\lambda_{x_{(2)},y_{(1)}}
(S\boxtimes S)\chi^\boxtimes\bigl(\iota_1(x_{(1)}y_{(1)})\iota_2(x_{(2)}y_{(2)})\bigr)\\
&=\sum\lambda_{x_{(2)},y_{(1)}}\lambda_{x_{(1)}y_{(1)},x_{(2)}y_{(2)}}
(S\boxtimes S)\bigl(\iota_1( x_{(2)}y_{(2)})\iota_2(x_{(1)}y_{(1)})\bigr)\\
&=\sum\lambda_{x_{(2)},y_{(1)}}
\lambda_{x_{(1)}y_{(1)},x_{(2)}y_{(2)}}
\iota_1\bigl(S(x_{(2)}y_{(2)})\bigr)\iota_2\bigl(S(x_{(1)}y_{(1)})\bigr)\\
&=\sum\lambda_{x_{(2)},y_{(1)}}
\lambda_{x_{(1)}y_{(1)},x_{(2)}y_{(2)}}
\lambda_{x_{(2)},y_{(2)}}
\lambda_{x_{(1)},y_{(1)}}
\iota_1\bigl(S(y_{(2)})S(x_{(2)})\bigr)\iota_2\bigl(S(y_{(1)})S(x_{(1)})\bigr).
\end{align*}
Now the right-hand sides of the two equations above coincide since (for every summand)
\begin{align*}
\deg(x)\deg(y)&+\deg(y_{(1)})\deg(y_{(2)})+
\deg(x_{(1)})\deg(x_{(2)})+\deg(x_{(2)})\deg(y_{(1)})\\
&=\deg(x_{(1)})\deg(y_{(1)})+\deg(x_{(1)})\deg(y_{(2)})\\
&\qquad+2\deg(x_{(2)})\deg(y_{(1)})+\deg(x_{(2)})\deg(y_{(2)})\\
&\qquad\quad+
\deg(y_{(1)})\deg(y_{(2)})+\deg(x_{(1)})\deg(x_{(2)})
\end{align*}
and
\begin{align*}
\deg(x_{(2)})\deg(y_{(1)})
&+\deg(x_{(1)}y_{(1)})\deg(x_{(2)}y_{(2)})
+\deg(x_{(2)})\deg(y_{(2)})+\deg(x_{(1)})\deg(y_{(1)})\\
&=2\deg(x_{(2)})\deg(y_{(1)})+\deg(x_{(1)})\deg(x_{(2)})+\deg(x_{(1)})\deg(y_{(2)})\\
&\qquad+
\deg(y_{(1)})\deg(y_{(2)})+\deg(x_{(2)})\deg(y_{(2)})+\deg(x_{(1)})\deg(y_{(1)}).
\end{align*}
Thus induction and Corollary \ref{cor-bazacA} can be used to easily finish the proof.
\end{proof}

\end{document}
